\section{家庭关系的法律适用}
\begin{enumerate}
  \item 婚姻。
        \begin{enumerate}
          \item 结婚、离婚。结婚手续任你意。诉讼离婚法院地，协离、财产自治地。共居共国行为地 。诉讼离婚只解决是否离婚的问题。
                \begin{kk}{效力的认定}
                  分形式要件和实质要件。看清题目问手续还是条件，还是混杂然后以偏概全！！
                  \al{定居甲国的华侨乙丙在甲国结婚，后来中国起诉离婚，认定其婚姻是否有效应适用甲国法 \cuo{没有明确区分是\xhx{手续还是条件}效力，不够准确}}。
                  \begin{enumerate}
                    \item 形式（手续）。一方当事人经居地、一方国籍国、婚姻缔结地均有效。
                    \item 实质（条件）。结婚时共居地 \yjt 共同国籍 \yjt 婚姻缔结地。按顺序。
                  \end{enumerate}
                \end{kk}
          \item \al{甲乙经居地均在中国，两人在丙国结婚。结婚条件应适用丙国法 \cuo}。
          \item 夫妻人身。共居共国行为地 \yjt 最密切联系原则。\al{改从夫姓、是否同居等}。
          \item 财产关系（动产、不动产）绝无可能用法院地法。
          \item 自治地。只能在经居地、国籍国、主要财产所在地选。
        \end{enumerate}
  \item 父母子女、监护和扶养（扩大解释）。父母监扶重弱者，父母先用共居地。共居地 \yjt 弱者经居地 \hz 国籍国。
  \item 涉外收养。
        \begin{enumerate}
          \item 程序。收养程序转亲书，条件手续重叠居。
                \begin{enumerate}
                  \item 通过外国收养组织向中国收养组织转交申请 + 提供家庭情况报告和证明。不允许直接来递材料。
                  \item 原则上共同亲自来华办理登记手续，可以委托但只能委托给收养关系的另一方配偶。
                  \item 应当与送养人订立书面收养协议。
                  \item 收养条件、程序。收养人和被收养人经居地。\al{久居上海的德国籍夫妇赴云南收养了女孩小梅}。
                \end{enumerate}
          \item 效力。收养效力收养居。收养时收养人经居地。
          \item 解除。解除被收法院地。收养时被收养人经居地 \hz 法院地。
        \end{enumerate}
  \item 继承。
        \begin{enumerate}
          \item 法定继承。法定不动所在地，动产死亡经居地。死亡时经常居所地。
          \item 遗嘱继承。遗嘱方式任你意，效力国籍或居地。立时亡时皆可用，管理绝产遗产地。
                \begin{kk}{效力的认定}
                  分形式要件和实质要件。现在很难依遗嘱方式去认定涉外遗嘱无效。
                  \begin{enumerate}
                    \item 形式。[立遗瞩、死亡时]，[经居地、国籍国] 排列组合、遗嘱行为地。
                    \item 实质。[立遗瞩、死亡时]，[经居地、国籍国] 排列组合。
                  \end{enumerate}
                  \tcblower
                  绝产。一个人死了遗产没人要。被继承人死亡时遗产所在地。
                \end{kk}
        \end{enumerate}
\end{enumerate}
